\title[mode=title]{The Threshold Does Not Transfer: Open-Set Test-Time Adaptation
Across Buildings}

\begin{abstract}
Deploying an activity recogniser in a new building means facing two shifts at once: the
sensors are in different places, and the occupant may do things the model was never
trained on. Open-set and universal domain adaptation address the second, and test-time
adaptation the first, but the methods for both are evaluated offline, with repeated passes
over a target set that is also the evaluation set. We study the combined shift under
constraints a deployment can meet: a source model trained once and frozen, a single online
pass over the unlabelled target stream in recording order, no source data, and no target
labels anywhere, including for the threshold that decides what counts as unknown.

Measured against a model trained on the target building at matched sample size, the
cross-building penalty on known activities is $1.6$ macro-F\textsubscript{1} points
entering one apartment and $36.0$ entering the other, a wider asymmetry than the
$0.6$ and $21.3$ of the closed-set case \cite{kbaier_closedset}. Four test-time methods
follow it: none improves on a frozen source model where there is no shift to recover.
One appears to, and the appearance is instructive. PETAL records the best open-set
H-score in one direction while classifying known activities roughly half as well, because
degrading a model uniformly lowers its confidence everywhere, pushes more windows below a
fixed threshold, and raises the unknown-detection term that H is bottlenecked by. The
standard open-set score is therefore maximised in part by damage, and pairing it with
known-class accuracy does not remove the problem; it only makes it visible.

The rejection rule fails for a different reason. Maximum softmax, entropy, maximum logit
and free energy are all at or below chance as novelty detectors here, with bootstrap
intervals excluding $0.5$ for three of the four: the model is as confident on activities it
never saw as on those it was trained on. A Mahalanobis distance to source class prototypes
is the only score above chance, and at a matched rejection rate it catches $1.58$ and
$1.17$ times as many unknowns as rejecting at random --- real, and far smaller than the
H-score comparison suggests, because two scores held at the same nominal percentile reject
very different fractions of a target stream.

That last observation is the main result. A threshold set to reject $1\%$ of source
validation rejects between $7\%$ and $63\%$ of the target stream, depending only on which
score carries it. Improving detection does not fix it: a source-only weighting rule that
gains $0.064$ AUROC in the direction with headroom \emph{loses} $0.11$ H at the threshold a
deployment would use, while the variant with the best H in both directions has
significantly worse AUROC in one. Detection quality and deployed open-set performance come
apart, so cross-domain threshold calibration is the binding constraint rather than one
difficulty among several. We recommend reporting the rejection rate, which is observable
without labels, beside every open-set score.
\end{abstract}

% ---------------------------------------------------------------------------
